\scanpage{77}
On est encore ramené à prouver ce dernier énoncé. On procède comme dans 2°, en introduisant la sous-cat. pleine $I'$ de $R$ telle que $\mathrm{Ob}\,I'=I$, d'où $T_R\longrightarrow T_{\widetilde I'}$ \uline{pl. fidèle}. Regardons $\widetilde I'\longrightarrow R$, il est \inserted{ess.}\deleted{surjectif sur les objets et} surjectif sur les Hom, \deleted{\unsure{donc}} est-il ess. surjectif ? OK. Si \deleted{le foncteur} $d'$ \deleted{\unsure{contient}} satisfait aux conditions écrites en marge plus haut, car alors on trouve un \inserted{$\Delta$-}foncteur pl. fidèle $\underbrace{\widetilde I'\Sigma^{-1}}_{\substack{\text{catégorie}\\\text{de fractions}\\\text{habituelle}}}\longrightarrow R$, \deleted{\unsure{dont}} qui doit donc être ess. surjectif. Donc on trouve

\diagram{$\begin{tikzcd}[ampersand replacement=\&,column sep=small,row sep=small]
\& \widetilde I' \arrow[ddl] \arrow[dl,dash] \arrow[dr,dash] \& \deleted{\unsure{I}}\\
I' \arrow[rr] \arrow[d,hook,"\text{pl. fid}"'] \& \& I' \arrow[d,"\substack{\text{ess. surj.}\\\text{et surjectif sur Hom}}"]\\
R \& \unsure{\subset|\longleftarrow|\hookleftarrow} \& \rho(R)
\end{tikzcd}$}

$I'$ coengendre strict\textsuperscript{mt} $R$ qui coengendre str. $\rho(R)$ donc $I'$ coengendre str. $\rho(R)$ (?)

\uline{Corollaire} \inserted{Si $d'$ satisfait aux conditions envisagées (p. ex. $d'=$ ens. des diagrammes finis)} On peut construire les théories algébriques pour $\Delta$ en trois pas, comme pour $\Delta_1$ : on prend une catégorie $I'$ ayant $I$ comme ens. d'objets (par 1° et 2°) et prend une catégorie de fractions de $\widetilde I'$.

% Two vertical strokes beside this paragraph.
Mais la m\textsuperscript{me} conclusion est valable si p. ex. $d'$ est formé de diagrammes discrets, stable par passage à un sous-diagramme \unsure{(|\{} \deleted{\unsure{diagrammes}}

Néanmoins, dans le cas d'un $d'$ général (m\textsuperscript{me} formé de diagrammes finis) il semble qu'il faille une construction récurrente transfinie

\deleted{\unsure{Th|Pb} \unsure{Question} \inserted{\uline{Question}} Supposons $d'$ petit. Alors une $R$ \deleted{telle que $R$} admettant une petite ss-catégorie $d'$-cogénératrice (donc \deleted{2-rep} correspondant à une théorie algébrique pour $\Delta$) est petite, car on \unsure{a}\ldots}
% archive sheet number: 76
